%% conclusion.tex
%%

%% ==================
\section{Conclusion and Future Work}
\label{sec:Conclusion}
%% ==================

This project asked whether a proactive Kubernetes autoscaler could keep most of its SLA-violation advantage over reactive HPA while substantially reducing the instability that Wanigasooriya \& Ekanayake's (2026) NimbusGuard exhibits and names but does not fix. Across 5 seeds, a stability-aware controller (exponential smoothing + hysteresis/cooldown, applied to the same forecasting model as an unfiltered baseline) cut scaling events by 57\% and pod-count volatility by 32\%, while still beating reactive HPA on SLA violations -- though giving back roughly half of the SLA-violation advantage the unfiltered proactive policy achieves.

The honest finding is a trade-off surface, not a single winning configuration: the smoothing factor and hysteresis/cooldown thresholds used here (0.4, 1 pod, 2 steps) are one point on a continuum between "aggressive and unstable" and "calm but slower to react," and were not exhaustively tuned.

\textbf{Future work:}
\begin{itemize}
  \item Sweep the smoothing/hysteresis/cooldown parameters to map the full SLA-vs-stability trade-off curve, rather than reporting one configuration.
  \item Evaluate on a real Kubernetes testbed with real workload traces, as NimbusGuard itself does, rather than a synthetic generator.
  \item Explore whether the cooldown/hysteresis rule could itself be learned (e.g. as a second, small model deciding \emph{when} to trust the predictor) rather than fixed, closing the gap to NimbusGuard's own RL-based approach while keeping the simplicity of not requiring a full DQN pipeline.
\end{itemize}
